\section{Exámenes y quizzes}
Tanto los exámenes como los quizzes deben realizarse a mano, de forma individual y sin consultar ningún material distinto al provisto por el profesor.

Puede completar el examen con lápiz, pero únicamente se aceptarán reclamos sobre respuestas escritas en tinta indeleble.

Se eximen de estas políticas los estudiantes que sean sujetos de adaptaciones autorizadas por la Universidad.

\subsection{Inicio y fin de exámenes}
Si un estudiante abandona el aula se da por terminado su examen; la única excepción es en caso de evacuación masiva a causa de una emergencia.

Los estudiantes pueden ingresar a presentar el examen en cualquier momento posterior a la hora de inicio del mismo, siempre y cuando ningún otro estudiante haya abandonado el aula. Tan pronto como un estudiante culmine su examen, ningún otro podrá ingresar al aula.

\subsection{Uso de lenguajes}
Se calificará el uso correcto de los lenguajes, tanto naturales como de programación, en las actividades evaluativas.

Como lenguaje natural se acepta el uso de español o inglés. Como lenguaje de programación se permite el uso de C++, Java, Python, JavaScript, TypeScript o GCL. La elección podrá variar entre apartados de una misma actividad evaluativa.

Cada apartado deberá redactarse íntegramente en un único lenguaje natural y, cuando corresponda, en un único lenguaje de programación. No es válido combinar idiomas naturales ni lenguajes de programación dentro de un mismo apartado.

\subsection{Política de ignorancia honesta}
Durante la presentación de un examen suele pasar que, tras leer el enunciado de una pregunta o problema, el estudiante es consciente de si su conocimiento le permite formular al menos una vaga aproximación a la respuesta, o de si por el contrario no tiene idea de cómo abordar el punto.

Ser capaz de reconocer los vacíos en el conocimiento propio es una habilidad sumamente útil en diversas facetas de la vida. A causa de eso, los exámenes evalúan y premian esa competencia. En caso de que considere que no dispone de la formación necesaria para dar respuesta a un apartado del examen, escribir explícitamente que no sabe y no responde le acreditará automáticamente el 20\% del valor del apartado.

Un ejemplo de respuesta que acredita esa valoración es: ``No sé''.

Por el contrario, no diligenciar el apartado o diligenciarlo con información incorrecta no otorgará ningún punto. Más aún, diligenciarlo con contenido incoherente o ininteligible podrá acarrear una penalización. Una respuesta manifiestamente absurda podrá recibir hasta el valor total del apartado en puntos \emph{negativos}.

\subsection{Política de información imprevista}
Ningún examen constituirá una sorpresa. Todo estudiante sabrá con el mayor de los detalles el temario a evaluar en cada una de las evaluaciones del curso.

En caso de que algún apartado de un examen requiera conocimiento que no se informó de antemano, cualquier estudiante podrá señalarlo en nombre de toda la sección. En caso de tener razón, el apartado transgresor será declarado inválido y la actividad evaluativa se calificará con base al resto de numerales.
